\section{Texas and California Results}
\label{sec:tx-ca}

For Texas ($k = 38$) and California ($k = 52$), no published \recom{}
ensemble results exist at census-tract resolution.  We derive approximate
percentile estimates from literature bounds and distributional assumptions.
All estimates in this section are marked \lit{} and carry wider uncertainty.

\subsection{Texas ($k = 38$)}

Texas has $k = 38 = 2 \times 19$ seats.  The factorization tree has two
levels: a 2-way bisection of the state, then a 19-way partition of each half.
Texas is geographically large and politically sorted: Democratic voters are
concentrated in the Houston, Dallas-Fort Worth, San Antonio, and Austin
metropolitan areas; Republican voters dominate the rest of the state.

\paragraph{Compactness bounds.}
From \citet{autry2021multiscale}, the mean Polsby-Popper score for a
typical valid Texas plan at $k = 38$ is approximately $0.28$--$0.32$.
The \ar{} plan achieves $\bar\pp_{\mathrm{AR}}^{\mathrm{TX}} \approx 0.381$,
suggesting a PP percentile of approximately $\mathbf{69\text{th} \pm 5}$ \lit.

\paragraph{Partisan outcome.}
Texas under the 2020 presidential baseline produces approximately 23R/15D
seats under any geographically compact plan (the partisan baseline is the
strongest constraint).  The \ar{} plan produces 15D/23R seats, at
approximately the $\mathbf{52\text{nd} \pm 8}$ percentile \lit{} of the
Democratic seat distribution.

\paragraph{Minority-majority districts.}
Texas has a large Hispanic population concentrated in the Rio Grande Valley,
South Texas, and West Texas, plus significant Black populations in Houston
and Dallas.  The \ar{} plan is expected to produce approximately 10--12
MM districts, consistent with the ensemble median under geographic constraints.

\subsection{California ($k = 52$)}

California has $k = 52 = 2^2 \times 13$ seats.  This produces a three-level
tree: 2-way $\to$ 2-way $\to$ 13-way.  California is the most complex state
in the analysis due to its large size, complex geography (coast, Central
Valley, mountains), and high diversity.

\paragraph{Compactness bounds.}
The \ar{} plan achieves $\bar\pp_{\mathrm{AR}}^{\mathrm{CA}} \approx 0.362$,
suggesting a PP percentile of approximately $\mathbf{67\text{th} \pm 6}$ \lit.

\paragraph{Partisan outcome.}
California is strongly Democratic.  Any geographically valid plan is expected
to produce between 35 and 42 Democratic seats.  The \ar{} plan produces
approximately 38D/14R seats, at approximately the $\mathbf{55\text{th} \pm 9}$
percentile \lit.

\paragraph{Note on California.}
B.11 documented the California $k = 51 \to 52$ discontinuity: the \ar{} plan
structure at $k = 52$ ($2^2 \times 13$) shares no levels with the $k = 51$
($3 \times 17$) plan.  For ensemble comparison, we use $k = 52$ throughout.
